\chapter{Implémentation et Démonstration}

\section*{Introduction du chapitre}
Ce dernier chapitre concrétise la conception en présentant la réalisation effective de
QuickTriage. Nous décrivons d'abord l'environnement et la pile technologique de développement,
puis l'implémentation du pipeline d'apprentissage automatique (préparation des données,
optimisation, calibration, résultats). Nous détaillons ensuite l'intégration du modèle dans le
backend J2EE à travers ses composants Java et l'échange REST. Le chapitre se termine par une
démonstration de la solution et une synthèse des artefacts produits.

\section{Environnement et pile technologique}

\subsection{Outils et versions}
La réalisation de QuickTriage a mobilisé un ensemble cohérent de technologies, réparties sur les
trois couches. Le tableau~\ref{tab:stack} en récapitule la pile, avec les versions retenues.

\begin{table}[h]
\centering
\renewcommand{\arraystretch}{1.35}
\begin{tabular}{p{3cm} p{5cm} p{5.5cm}}
\toprule
\textbf{Couche} & \textbf{Technologie} & \textbf{Rôle} \\
\midrule
\multirow{6}{*}{Apprentissage}
 & Python 3 & Langage du pipeline ML \\
 & pandas, NumPy & Manipulation des données \\
 & scikit-learn & Prétraitement, métriques, SMOTE (\textit{imbalanced-learn}) \\
 & XGBoost & Modèle de classification \\
 & Optuna & Optimisation des hyperparamètres \\
 & ONNX / skl2onnx & Export du modèle interopérable \\
\midrule
\multirow{4}{*}{Backend}
 & Jakarta EE & Plateforme d'entreprise \\
 & WildFly 10 & Serveur d'application \\
 & JAX-RS / EJB & API REST et logique métier \\
 & ONNX Runtime 1.18.0 & Moteur d'inférence (\texttt{com.microsoft.onnxruntime}) \\
\midrule
\multirow{2}{*}{Frontend}
 & Android (Android Studio) & Application mobile (Java/Kotlin) \\
 & Retrofit / OkHttp & Consommation de l'API REST \\
\midrule
Outillage & Maven, Git & Construction et gestion de versions \\
\bottomrule
\end{tabular}
\caption{Pile technologique de QuickTriage.}
\label{tab:stack}
\end{table}

\section{Implémentation du pipeline d'apprentissage automatique}

\subsection{Préparation et augmentation des données}
Le jeu de données \textit{KTAS Emergency Triage} comporte initialement 1~267 patients réels,
fortement déséquilibrés entre les classes. Pour corriger ce déséquilibre, une stratégie
d'augmentation hybride a été mise en œuvre, combinant la génération de patients synthétiques
calibrés sur les distributions réelles et l'algorithme SMOTE. Le résultat est un jeu équilibré de
3~000 patients, soit 600 par classe métier. Le listing~\ref{lst:smote} illustre le principe du
rééquilibrage par SMOTE avec un nombre de voisins adaptatif.

\begin{listing}[H]
\begin{minted}[bgcolor=bg,fontsize=\small,linenos]{python}
from imblearn.over_sampling import SMOTE
from collections import Counter

# Determination d'un k_neighbors adaptatif (robuste aux petites classes)
counts = Counter(y_train)
min_class_size = min(counts.values())
k_neighbors = max(1, min(5, min_class_size - 1))

smote = SMOTE(random_state=42, k_neighbors=k_neighbors)
X_balanced, y_balanced = smote.fit_resample(X_train, y_train)

print("Distribution apres reequilibrage :", Counter(y_balanced))
# -> {CRITIQUE: 600, URGENT: 600, NON_URGENT: 600}
\end{minted}
\caption{Rééquilibrage des classes par SMOTE avec \texttt{k\_neighbors} adaptatif.}
\label{lst:smote}
\end{listing}

\subsection{Optimisation des hyperparamètres avec Optuna}
Le modèle XGBoost a été optimisé par Optuna sur \textbf{80 essais}. La fonction objectif maximise
le F1-score pondéré obtenu en validation croisée. Le listing~\ref{lst:optuna} en présente une forme
simplifiée.

\begin{listing}[H]
\begin{minted}[bgcolor=bg,fontsize=\small,linenos]{python}
import optuna
from xgboost import XGBClassifier
from sklearn.model_selection import cross_val_score

def objective(trial):
    params = {
        "n_estimators":  trial.suggest_int("n_estimators", 100, 600),
        "max_depth":     trial.suggest_int("max_depth", 3, 10),
        "learning_rate": trial.suggest_float("learning_rate", 1e-3, 0.3, log=True),
        "subsample":     trial.suggest_float("subsample", 0.6, 1.0),
        "colsample_bytree": trial.suggest_float("colsample_bytree", 0.6, 1.0),
    }
    model = XGBClassifier(**params, eval_metric="mlogloss", random_state=42)
    score = cross_val_score(model, X_balanced, y_balanced,
                            cv=5, scoring="f1_weighted").mean()
    return score

study = optuna.create_study(direction="maximize")
study.optimize(objective, n_trials=80)
print("Meilleur essai :", study.best_trial.number)   # -> 64
print("Meilleurs parametres :", study.best_params)
\end{minted}
\caption{Optimisation des hyperparamètres de XGBoost avec Optuna (80 essais).}
\label{lst:optuna}
\end{listing}

\noindent Le meilleur essai (\textbf{n\textsuperscript{o}~64}) retient les hyperparamètres
suivants : \texttt{n\_estimators}~=~458, \texttt{max\_depth}~=~7 et
\texttt{learning\_rate}~$\approx$~0,0113. Ces valeurs traduisent un modèle relativement profond et
nombreux en arbres, mais à pas d'apprentissage faible, gage de stabilité.

\subsection{Calibration du seuil orientée sécurité patient}
Conformément au parti pris de minimiser le sous-triage, le seuil de décision de la classe
\textsc{Critique} a été abaissé à \textbf{0,30} au lieu de la règle usuelle de l'\textit{argmax}.
Concrètement, dès que la probabilité estimée d'appartenance à la classe \textsc{Critique} atteint
0,30, le patient est classé comme critique. Ce réglage privilégie délibérément le rappel des cas
graves au prix d'un léger sur-triage, jugé acceptable d'un point de vue clinique.

\subsection{Résultats et évaluation}
Le tableau~\ref{tab:metrics} synthétise les performances du modèle final. L'évaluation met
l'accent sur la classe \textsc{Critique}, dont le rappel constitue l'indicateur de sécurité
déterminant.

\begin{table}[h]
\centering
\renewcommand{\arraystretch}{1.35}
\begin{tabular}{p{6cm} c}
\toprule
\textbf{Métrique} & \textbf{Valeur} \\
\midrule
Exactitude globale (\textit{accuracy}) & 73,8~\% \\
F1-score pondéré (\textit{weighted}) & 0,729 \\
Rappel classe \textsc{Critique} & 83,9~\% \\
F1-score classe \textsc{Critique} & 0,80 \\
Seuil de décision \textsc{Critique} & 0,30 \\
\bottomrule
\end{tabular}
\caption{Performances du modèle XGBoost final de QuickTriage.}
\label{tab:metrics}
\end{table}

\begin{successbox}
Avec un \textbf{rappel de 83,9~\%} sur la classe \textsc{Critique}, le modèle identifie
correctement la grande majorité des patients en danger vital, ce qui constitue l'objectif
prioritaire d'un système d'aide au triage.
\end{successbox}

\noindent La matrice de confusion (figure~\ref{fig:confusion}) détaille la répartition des
prédictions par classe et confirme la rareté des erreurs de sous-triage les plus graves
(patients critiques classés non urgents).

\figplaceholder{[Figure 6.1 -- Matrice de confusion du modèle final sur le jeu de test
(classes CRITIQUE, URGENT, NON\_URGENT). À insérer : image exportée depuis le notebook
d'évaluation.]}
\begin{figure}[h]
\centering
\caption{Matrice de confusion du modèle final.}
\label{fig:confusion}
\end{figure}

\subsection{Export ONNX}
Une fois le modèle entraîné et validé, il est exporté au format ONNX afin d'être consommé par le
backend Java. Le listing~\ref{lst:onnx} montre l'export via \texttt{skl2onnx}. L'artefact produit,
\texttt{triage\_final\_v3.onnx}, occupe environ 139~Ko.

\begin{listing}[H]
\begin{minted}[bgcolor=bg,fontsize=\small,linenos]{python}
from skl2onnx import convert_sklearn
from skl2onnx.common.data_types import FloatTensorType

initial_type = [("float_input", FloatTensorType([None, 23]))]
onnx_model = convert_sklearn(best_model, initial_types=initial_type,
                             options={type(best_model): {"zipmap": False}})

with open("triage_final_v3.onnx", "wb") as f:
    f.write(onnx_model.SerializeToString())
\end{minted}
\caption{Export du modèle au format ONNX (entrée \texttt{float\_input} de dimension 23).}
\label{lst:onnx}
\end{listing}

\section{Intégration du modèle dans le backend J2EE}

\subsection{Chargement du modèle : \texttt{OnnxModelLoader}}
Le bean \texttt{OnnxModelLoader} charge le modèle une seule fois au démarrage du serveur grâce aux
annotations \texttt{@Singleton} et \texttt{@Startup}. La session ONNX est ensuite réutilisée pour
toutes les inférences (listing~\ref{lst:loader}).

\begin{listing}[H]
\begin{minted}[bgcolor=bg,fontsize=\small,linenos,breaklines]{java}
@Singleton
@Startup
public class OnnxModelLoader {

    private OrtEnvironment env;
    private OrtSession session;

    @PostConstruct
    public void init() {
        try {
            env = OrtEnvironment.getEnvironment();
            byte[] model = getClass().getResourceAsStream("/triage_final_v3.onnx")
                                     .readAllBytes();
            session = env.createSession(model, new OrtSession.SessionOptions());
        } catch (Exception e) {
            throw new IllegalStateException("Echec du chargement du modele ONNX", e);
        }
    }

    public OrtSession getSession() {
        return session;
    }

    @PreDestroy
    public void close() throws Exception {
        if (session != null) session.close();
        if (env != null) env.close();
    }
}
\end{minted}
\caption{Chargement unique du modèle ONNX au démarrage (\texttt{OnnxModelLoader}).}
\label{lst:loader}
\end{listing}

\subsection{Logique métier : \texttt{TriagePredictionService}}
L'EJB \texttt{@Stateless} \texttt{TriagePredictionService} construit le vecteur de
caractéristiques, applique le prétraitement issu des métadonnées, exécute l'inférence et applique
le seuil de criticité (listing~\ref{lst:service}).

\begin{listing}[H]
\begin{minted}[bgcolor=bg,fontsize=\small,linenos,breaklines]{java}
@Stateless
public class TriagePredictionService {

    @Inject
    private OnnxModelLoader modelLoader;

    private final ModelMetadata meta = ModelMetadata.load("/metadata_v3.json");

    public TriageResponse predict(TriageRequest req) throws Exception {
        float[] features = normalize(buildFeatures(req));

        OrtSession session = modelLoader.getSession();
        OrtEnvironment env = OrtEnvironment.getEnvironment();
        OnnxTensor input = OnnxTensor.createTensor(env, new float[][]{features});

        try (OrtSession.Result result =
                 session.run(Collections.singletonMap("float_input", input))) {
            float[] proba = extractProbabilities(result);
            return applyThreshold(proba);
        }
    }

    private TriageResponse applyThreshold(float[] proba) {
        TriageResponse resp = new TriageResponse();
        if (proba[KtasClass.CRITIQUE.ordinal()] >= meta.getSeuilCritique()) {
            resp.setKtasClass(KtasClass.CRITIQUE);
            resp.setCritical(true);
        } else {
            resp.setKtasClass(KtasClass.values()[argmax(proba)]);
        }
        resp.setProbabilities(toMap(proba));
        resp.setModelVersion(meta.getModelVersion());
        return resp;
    }
}
\end{minted}
\caption{Logique métier de prédiction et application du seuil (\texttt{TriagePredictionService}).}
\label{lst:service}
\end{listing}

\subsection{Exposition REST : \texttt{TriageResource}}
La ressource JAX-RS \texttt{TriageResource} expose le point d'entrée \texttt{POST
/api/triage/predict} et délègue le traitement au service (listing~\ref{lst:resource}).

\begin{listing}[H]
\begin{minted}[bgcolor=bg,fontsize=\small,linenos,breaklines]{java}
@Path("/api/triage")
@Produces(MediaType.APPLICATION_JSON)
@Consumes(MediaType.APPLICATION_JSON)
public class TriageResource {

    @Inject
    private TriagePredictionService predictionService;

    @POST
    @Path("/predict")
    public Response predict(TriageRequest request) {
        try {
            TriageResponse response = predictionService.predict(request);
            return Response.ok(response).build();
        } catch (Exception e) {
            return Response.status(Response.Status.INTERNAL_SERVER_ERROR)
                           .entity("Erreur lors de la prediction").build();
        }
    }
}
\end{minted}
\caption{Point d'entrée REST de prédiction (\texttt{TriageResource}).}
\label{lst:resource}
\end{listing}

\subsection{Dépendance Maven}
L'intégration du moteur d'inférence repose sur la dépendance Maven présentée au
listing~\ref{lst:maven}.

\begin{listing}[H]
\begin{minted}[bgcolor=bg,fontsize=\small]{xml}
<dependency>
    <groupId>com.microsoft.onnxruntime</groupId>
    <artifactId>onnxruntime</artifactId>
    <version>1.18.0</version>
</dependency>
\end{minted}
\caption{Dépendance Maven pour ONNX Runtime.}
\label{lst:maven}
\end{listing}

\subsection{Format d'échange REST}
L'échange entre le client Android et le backend repose sur JSON. Le listing~\ref{lst:json}
présente un exemple de requête et la réponse correspondante pour un patient critique.

\begin{listing}[H]
\begin{minted}[bgcolor=bg,fontsize=\small]{json}
// Requete : POST /api/triage/predict
{
  "age": 72, "mental": 3, "nrsPain": 8,
  "sbp": 85, "dbp": 50, "hr": 124, "rr": 28,
  "bt": 38.9, "saturation": 90, "arrivalMode": 3
}

// Reponse : 200 OK
{
  "ktasClass": "CRITIQUE",
  "ktasLevels": "1-2",
  "delaiMax": 15,
  "confidence": 0.84,
  "isCritical": true,
  "probabilities": { "CRITIQUE": 0.84, "URGENT": 0.12, "NON_URGENT": 0.04 },
  "modelVersion": "v3"
}
\end{minted}
\caption{Exemple d'échange REST JSON (requête et réponse).}
\label{lst:json}
\end{listing}

\section{Synthèse des artefacts produits}
Le tableau~\ref{tab:artefacts} récapitule les principaux artefacts de la chaîne QuickTriage.

\begin{table}[h]
\centering
\renewcommand{\arraystretch}{1.35}
\begin{tabular}{p{4.5cm} p{3cm} p{6cm}}
\toprule
\textbf{Artefact} & \textbf{Type} & \textbf{Description} \\
\midrule
\texttt{triage\_final\_v3.onnx} & Modèle (139 Ko) & Modèle XGBoost exporté, consommé par le backend \\
\texttt{metadata\_v3.json} & Métadonnées & Scaler, médianes d'imputation, seuil, version \\
\texttt{OnnxModelLoader.java} & Backend & Chargement unique du modèle (Singleton) \\
\texttt{TriagePredictionService.java} & Backend & Logique métier d'inférence (EJB) \\
\texttt{TriageResource.java} & Backend & Point d'entrée REST (JAX-RS) \\
Application Android & Frontend & Saisie et restitution mobile \\
\bottomrule
\end{tabular}
\caption{Principaux artefacts de la solution QuickTriage.}
\label{tab:artefacts}
\end{table}

\section{Démonstration}
Cette section illustre le fonctionnement de l'application mobile sur un parcours type. Les captures
d'écran (à insérer) présentent successivement l'écran de saisie des constantes, l'écran de
résultat pour un cas critique et l'écran de résultat pour un cas non urgent.

\figplaceholder{[Figure 6.2 -- Écran de saisie des constantes vitales dans l'application Android.]}
\begin{figure}[h]
\centering
\caption{Application QuickTriage -- écran de saisie des constantes.}
\label{fig:demo-saisie}
\end{figure}

\figplaceholder{[Figure 6.3 -- Écran de résultat pour un patient CRITIQUE (bandeau rouge,
alerte de criticité, délai immédiat, confiance et probabilités).]}
\begin{figure}[h]
\centering
\caption{Application QuickTriage -- résultat pour un cas critique.}
\label{fig:demo-critique}
\end{figure}

\figplaceholder{[Figure 6.4 -- Écran de résultat pour un patient NON\_URGENT (bandeau vert).]}
\begin{figure}[h]
\centering
\caption{Application QuickTriage -- résultat pour un cas non urgent.}
\label{fig:demo-nonurgent}
\end{figure}

\noindent La démonstration confirme le comportement attendu de bout en bout : la saisie des
constantes déclenche un appel à l'API, l'inférence est restituée en une fraction de seconde et le
code couleur permet une lecture immédiate de la priorité. Le cas critique déclenche bien l'alerte
\texttt{isCritical}, conformément à la règle de seuil calibrée.

\section*{Conclusion du chapitre}
Ce chapitre a présenté la réalisation concrète de QuickTriage : la pile technologique,
l'implémentation du pipeline d'apprentissage (augmentation hybride, optimisation Optuna,
calibration orientée sécurité, export ONNX), l'intégration du modèle dans le backend J2EE via ses
trois composants Java, l'échange REST JSON, ainsi qu'une démonstration de l'application mobile. La
solution couvre ainsi l'intégralité de la chaîne, de la donnée brute jusqu'à l'aide à la décision
restituée au soignant. La conclusion générale dresse le bilan de ce travail et en esquisse les
perspectives.
